% !TeX TXS-program:compile = txs:///pdflatex/[--shell-escape]

\documentclass[11pt, letterpaper]{article}

\usepackage{minted}
\usepackage[utf8]{inputenc}
\usepackage[T1]{fontenc}
\usepackage{lmodern}
\usepackage{graphicx}
\usepackage{longtable}
\usepackage{wrapfig}
\usepackage{rotating}
\usepackage{amsmath}
\usepackage{textcomp}
\usepackage{amssymb}
\usepackage{hyperref}
\usepackage[round]{natbib}
\usepackage{subcaption}


\title{\bfseries Tarea}
\author{Ángel García Báez}
\date{\today}
\setcounter{tocdepth}{3} 

\begin{document}
	
\textbf{Cruza Intermedia}

Sean:
\[
P_1 = (v_1, v_2, \ldots, v_m)
\]
\[
P_2 = (w_1, w_2, \ldots, w_m)
\]

Si se cruzan en la posición \( k \), los hijos son:

\[
H_1 = (v_1, \ldots, v_k, \, a w_{k+1} + (1-a) v_{k+1}, \ldots, a w_m + (1-a) v_m), \quad a \in [0, 1]
\]

\[
H_2 = (w_1, \ldots, w_k, \, a v_{k+1} + (1-a) w_{k+1}, \ldots, a v_m + (1-a) w_m), \quad a \in [0, 1]
\]

\textbf{Cruza Aritmética Simple}

El valor de \( a \) depende de la comparación entre \( v_k \) y \( w_k \):

\[
a \in
\begin{cases}
	\left[ \max(\alpha, \beta), \min(\gamma, \delta) \right], & \text{si } v_k > w_k \\
	[0, 0], & \text{si } v_k = w_k \\
	\left[ \max(\gamma, \delta), \min(\alpha, \beta) \right], & \text{si } v_k < w_k
\end{cases}
\]

\text{Donde:}
\[
\alpha = \frac{l_k^w - w_k}{v_k - w_k}, \quad
\beta = \frac{u_k^v - v_k}{w_k - v_k}, \quad
\gamma = \frac{l_k^v - v_k}{w_k - v_k}, \quad
\delta = \frac{u_k^w - w_k}{v_k - w_k}
\]

\text{Rangos de validez:}
\[
v_k \in [l_k^v, u_k^v], \quad
w_k \in [l_k^w, u_k^w]
\]


\textbf{Mutación Simple}

Sea un individuo \( x = (x_1, x_2, \ldots, x_m) \), se selecciona un gen \( x_i \) y se modifica de la siguiente forma:

\[
x_i' =  \Delta
\]

donde \( \Delta \) es un pequeño valor aleatorio, típicamente:

\[
\Delta \sim \mathcal{U}(Lower, Upper)
\]

o

\[
\Delta \sim \mathcal{N}(0, \sigma^2)
\]

$\text{Rango de validez:} \quad x_i' \in [l_i, u_i]$


\textbf{Mutación por Reordenamiento}

Sea un individuo representado por una permutación:

\[
x = (x_1, x_2, \ldots, x_i, \ldots, x_j, \ldots, x_m)
\]

Se selecciona aleatoriamente un subconjunto de posiciones \( i \) a \( j \), donde \( 1 \leq i < j \leq m \), y se reordena esa subsecuencia:

\[
x' = (x_1, x_2, \ldots, \underbrace{\text{perm}(x_i, \ldots, x_j)}_{\text{reordenamiento}}, \ldots, x_m)
\]

$\text{Donde } \text{perm}(x_i, \ldots, x_j) \text{ es una permutación aleatoria de los elementos } x_i, \ldots, x_j.$



\section*{Escalamiento Sigma}

\[
v_{ei} =
\begin{cases}
	1 + \dfrac{\text{Apt}_i - \overline{\text{Apt}}}{2\sigma}, & \text{si } \sigma^2 \neq 0 \\
	1, & \text{si } \sigma = 0
\end{cases}
\]

\[
\sigma = \sqrt{ \dfrac{n \sum \text{Apt}_i^2 - \left( \sum \text{Apt}_i \right)^2}{n^2} }
\]

\section*{Selección de Jerarquías}

\textbf{Algoritmo:}
\begin{enumerate}
	\item Ordenar (jerarquizar) la población con base en su aptitud, de 1 a \( N \) (donde 1 representa al menos apto).
	\item Elegir \( \text{Max} \) tal que \( 1 \leq \text{Max} \leq 2 \)
	\item Calcular \( \text{Min} = 2 - \text{Max} \)
	\item Calcular el valor esperado:
	\[
	\text{Valesp}_{(i,t)} = \text{Min} + (\text{Max} - \text{Min}) \cdot \dfrac{\text{jerarquía}_{(i,t)} - 1}{N - 1}
	\]
\end{enumerate}

Este método considera únicamente el orden, ignorando los valores reales de aptitud.

\section*{Selección de Boltzmann}

\[
\text{Valesp}_{(i,t)} = \dfrac{e^{f(i)/T}}{ \langle e^{f(i)/T} \rangle^t }
\]

Donde:
\begin{itemize}
	\item \( T \) es la temperatura.
	\item \( \langle \cdot \rangle^t \) denota el promedio de la población en la generación \( t \).
\end{itemize}

\section*{Selección por Ruleta}

\begin{enumerate}
	\item Calcular la suma de los valores esperados \( T \).
	\item Repetir \( N \) veces (donde \( N \) es el tamaño de la población):
	\begin{itemize}
		\item Generar un número aleatorio \( r \in [0, T] \).
		\item Recorrer la población sumando los valores esperados hasta que la suma sea mayor o igual a \( r \).
		\item El individuo que exceda este umbral es seleccionado.
	\end{itemize}
\end{enumerate}

\[
v_{ei} = \dfrac{\text{Aptitud}_i}{\overline{\text{Aptitud}}}, \quad \sum v_{ei} \approx \text{popsize}
\]

\section*{Selección por Torneo}

\begin{itemize}
	\item Seleccionar aleatoriamente \( t \) individuos.
	\item \textbf{Determinística:} Se elige al individuo con la mejor aptitud.
	\item \textbf{Probabilística:} El mejor individuo se selecciona con probabilidad \( p \), donde \( 0.5 \leq p < 1 \).
\end{itemize}

\section*{Sobrante Estocástico}

\[
\text{Valesp}_{(i)} = \dfrac{f_i}{f}
\]

\textbf{Algoritmo:}
\begin{enumerate}
	\item Asignar de manera determinística los valores enteros de los valores esperados.
	\item Los valores decimales restantes se usan probabilísticamente para llenar la población.
	\item \textbf{Sin reemplazo:} Usar un lanzamiento (flip) con las partes decimales.
	\item \textbf{Con reemplazo:} Construir una ruleta con las partes decimales y seleccionar los padres restantes.
\end{enumerate}

\section*{Selección Universal Estocástica}

\section*{Selección Universal Estocástica}

Sea \( N \) el número de individuos en la población y \( V_i \) el valor esperado del individuo \( i \). Primero calculamos:

\[
T = \sum_{i=1}^N V_i
\]

Luego, se generan \( N \) punteros equiespaciados a partir de un valor aleatorio \( r \in [0, \frac{T}{N}) \):

\[
\text{Punteros: } p_k = r + \dfrac{k \cdot T}{N}, \quad k = 0, 1, \ldots, N-1
\]

Después, para cada puntero \( p_k \), se selecciona el individuo \( i \) tal que:

\[
\sum_{j=1}^{i-1} V_j < p_k \leq \sum_{j=1}^{i} V_j
\]

Este individuo \( i \) es entonces seleccionado como parte de la nueva población.



\textbf{Algoritmo:}
\begin{verbatim}
	ptr = Rand(); /* Aleatorio en (0,1] */
	for (sum = 0, i = 1; i <= n; i++)
	for (sum += Valesp(i,t); sum > ptr; ptr++)
	Seleccionar(i);
\end{verbatim}

\section*{Muestreo Determinístico}

\textbf{Algoritmo:}
\begin{enumerate}
	\item Calcular:
	\[
	p_{\text{select}} = \dfrac{f_i}{\sum f}
	\]
	\item Calcular:
	\[
	\text{Valesp}_i = p_{\text{select}} \cdot n
	\]
	\item Asignar la parte entera de \( \text{Valesp}_i \) de forma determinística.
	\item Ordenar la población según la parte decimal de \( \text{Valesp}_i \) (de mayor a menor).
	\item Seleccionar los padres faltantes de la parte superior de la lista.
\end{enumerate}

\section*{Conversión de Binario a Real}

\subsection*{Discretización}

\[
L = \left\lfloor \log_2 \left[ (l_{\text{sup}} - l_{\text{inf}}) \cdot 10^{\text{precisión}} \right] + 0.9 \right\rfloor
\]

\[
X_{\text{real}} = l_{\text{inf}} + \dfrac{X_{\text{ent}} \cdot (l_{\text{sup}} - l_{\text{inf}})}{2^L - 1}
\]

	
	
\end{document}



